\documentclass[a4paper,12pt]{article}
\usepackage{indentfirst}
\usepackage[T1]{fontenc}

\begin{document}

\section{Conclus\~oes}


Este trabalho desenvolveu uma plataforma vest\'ivel modular multissensor baseada no XIAO ESP32-C6, no ESP-IDF e no FreeRTOS. O problema tratado foi a integra\c{c}\~ao de sensores heterog\^eneos, interface gr\'afica e servi\c{c}os compartilhados sem concentrar o processamento no n\'ucleo e sem interromper deliberadamente as fun\c{c}\~oes independentes quando um sensor n\~ao est\'a dispon\'ivel.


\subsection{Atendimento aos objetivos}


A vers\~ao implementada separa n\'ucleo, servi\c{c}os, drivers, processamento e telas. MAX30102, DS18B20, LTR390, b\'ussola e ped\^ometro seguem o ciclo comum de inicializa\c{c}\~ao, cria\c{c}\~ao e exibi\c{c}\~ao de tela, respeitadas as diferen\c{c}as de transporte. O despacho gr\'afico atua sobre a tela ativa, enquanto tarefas pr\'oprias executam as aquisi\c{c}\~oes. A navega\c{c}\~ao, as quatro p\'aginas de menu, o RTC e as telas funcionais foram confirmados pelo firmware e pelas capturas da interface.


O barramento I$^2$C foi centralizado em um servi\c{c}o comum com mutex, heran\c{c}a de prioridade, limites de espera, varredura no boot e recupera\c{c}\~ao p\'os-NACK. A an\'alise est\'atica confirmou esses caminhos, e os ensaios com remo\c{c}\~ao individual e combinada dos sensores mostraram que a aus\^encia de um ou mais componentes n\~ao impediu o uso do sistema principal nem dos m\'odulos independentes. Como n\~ao foram induzidas falhas durante a opera\c{c}\~ao nem executadas repeti\c{c}\~oes sob um protocolo estat\'istico, a evid\^encia sustenta a presen\c{c}a dos mecanismos e comprova a conten\c{c}\~ao funcional nas condi\c{c}\~oes testadas, mas n\~ao uma taxa de recupera\c{c}\~ao ou confiabilidade.


A extensibilidade e o acoplamento foram avaliados estruturalmente. No estudo retrospectivo do MAX30102, seis arquivos de processamento, com 678 linhas, foram reutilizados sem altera\c{c}\~ao; a integra\c{c}\~ao exigiu seis pontos expl\'icitos fora da pasta do m\'odulo e n\~ao modificou os demais m\'odulos sensoriais. Esse resultado demonstra localiza\c{c}\~ao das altera\c{c}\~oes no caso estudado, mas n\~ao mede o tempo necess\'ario para integrar um componente arbitr\'ario.


A compila\c{c}\~ao limpa produziu uma imagem grav\'avel de 784.240 bytes, correspondente a 74,79\% da parti\c{c}\~ao de 1 MiB. O resultado confirma o enquadramento global do firmware integrado. N\~ao foram obtidos incrementos por m\'odulo nem medidas compar\'aveis de CPU e heap por cen\'ario, de modo que a caracteriza\c{c}\~ao de recursos permanece parcial.


\subsection{Contribui\c{c}\~oes e alcance dos resultados}


A contribui\c{c}\~ao principal \'e a organiza\c{c}\~ao concreta de uma plataforma multissensor em contratos e servi\c{c}os compartilhados. A implementa\c{c}\~ao concentra em cada m\'odulo o acesso ao dispositivo, o processamento, o estado e a apresenta\c{c}\~ao, enquanto o n\'ucleo coordena inicializa\c{c}\~ao, registro e navega\c{c}\~ao. A placa prot\'otipo autoral e o inv\'olucro impresso em 3D tamb\'em foram realizados e integram o resultado de implementa\c{c}\~ao.


Os ensaios sensoriais demonstraram funcionamento, n\~ao valida\c{c}\~ao cl\'inica ou metrol\'ogica. O PPG apresentou componente puls\'atil coerente entre os canais e indica\c{c}\~oes pr\'oximas \`as do ox\'imetro no ensaio limitado; o DS18B20 respondeu aos est\'imulos t\'ermicos e diferiu $-0{,}9\,^\circ$C do comparador em um ponto; o LTR390 respondeu \`a luz e ao ultravioleta, com satura\c{c}\~ao da configura\c{c}\~ao em torno de 52 klux; a b\'ussola respondeu \`a mudan\c{c}a de dire\c{c}\~ao, com diferen\c{c}as de 16$^\circ$ e 22$^\circ$ nas duas orienta\c{c}\~oes ensaiadas; e o ped\^ometro registrou 131 e 133 passos para uma contagem manual de 140. O consumo total variou de 90 a 210 mA entre os nove cen\'arios observados, mas n\~ao foi realizada uma curva completa de descarga.


\subsection{Limita\c{c}\~oes}


As restri\c{c}\~oes de prazo delimitaram o escopo experimental desta vers\~ao. Permaneceram sem caracteriza\c{c}\~ao: CPU e heap por cen\'ario, mem\'oria incremental por m\'odulo, s\'erie controlada de tempo de boot, recupera\c{c}\~ao I$^2$C sob falhas induzidas, ensaio prolongado com roteiro completo e curva de descarga. A placa e o inv\'olucro n\~ao passaram por ensaios mec\^anicos ou ambientais, e o backlight permaneceu ligado.


As estimativas de frequ\^encia card\'iaca e SpO$_2$ n\~ao foram validadas clinicamente e permanecem suscet\'iveis a movimento, contato e condi\c{c}\~oes \'opticas. Temperatura, ilumin\^ancia, UVI e \emph{heading} n\~ao foram comparados com refer\^encias rastre\'aveis. A b\'ussola n\~ao compensa a inclina\c{c}\~ao nem aplica uma corre\c{c}\~ao magn\'etica matricial completa. O ped\^ometro foi avaliado com um participante, um ritmo e duas repeti\c{c}\~oes, usa par\^ametros fixos e perde a contagem ap\'os reinicializa\c{c}\~ao.


\subsection{Trabalhos futuros}


Os trabalhos futuros s\~ao priorizados pelas limita\c{c}\~oes observadas e por expans\~oes compat\'iveis com a arquitetura.


\subsubsection{Caracteriza\c{c}\~ao e robustez}


\begin{itemize}
    \item medir carga de CPU, heap livre e menor heap observado em cada cen\'ario funcional;
    \item quantificar o incremento de c\'odigo, dados e mem\'oria din\^amica associado a cada m\'odulo;
    \item induzir falhas I$^2$C controladas e medir detec\c{c}\~ao, recupera\c{c}\~ao, continuidade dos demais endere\c{c}os, resets e watchdogs;
    \item executar ensaio prolongado com roteiro repet\'ivel de navega\c{c}\~ao e m\'odulos ativos; e
    \item realizar uma curva completa de descarga em cen\'ario funcional definido.
\end{itemize}


\subsubsection{Energia}


\begin{itemize}
    \item avaliar o controle efetivo do backlight em uma revis\~ao de hardware compat\'ivel;
    \item implementar desligamento seletivo de sensores, modos de baixo consumo e despertar pelo RTC; e
    \item considerar um contador de coulombs para estimar a carga com base em corrente acumulada, comparando-o com a indica\c{c}\~ao baseada somente em tens\~ao.
\end{itemize}


\subsubsection{Armazenamento, conectividade e localiza\c{c}\~ao}


\begin{itemize}
    \item integrar o cart\~ao microSD para armazenar hist\'oricos, dados de ensaio e recursos gr\'aficos hoje mantidos na flash, quando o acesso externo for tecnicamente apropriado;
    \item coordenar o SPI compartilhado entre cart\~ao e display por sinais CS independentes e exclus\~ao m\'utua, e medir quanto da flash interna \'e efetivamente liberado; a migra\c{c}\~ao de arquivos n\~ao desloca automaticamente o c\'odigo execut\'avel;
    \item integrar Wi-Fi e Bluetooth do ESP32-C6 para transfer\^encia de dados, configura\c{c}\~ao e sincroniza\c{c}\~ao, avaliando consumo, seguran\c{c}a e impacto sobre a concorr\^encia; e
    \item investigar um m\'odulo GNSS com suporte a RTK para rastreamento de trajet\'orias. A avalia\c{c}\~ao deve considerar antena, interface, consumo, fonte das corre\c{c}\~oes, opera\c{c}\~ao em c\'eu aberto e compara\c{c}\~ao com refer\^encia apropriada, sem antecipar precis\~ao antes do ensaio.
\end{itemize}


\subsubsection{Algoritmos e evolu\c{c}\~ao modular}


\begin{itemize}
    \item implementar compensa\c{c}\~ao de inclina\c{c}\~ao da b\'ussola e avaliar calibra\c{c}\~ao magn\'etica matricial;
    \item persistir o ped\^ometro, reinici\'a-lo diariamente pelo RTC, investigar par\^ametros adaptativos e ampliar os ensaios para mais participantes, ritmos e posi\c{c}\~oes;
    \item acrescentar \'indice de qualidade do PPG, detec\c{c}\~ao de movimento e rejei\c{c}\~ao de artefatos; e
    \item avaliar um registro declarativo ou din\^amico de m\'odulos para reduzir os pontos de altera\c{c}\~ao no n\'ucleo, medindo o custo da abstra\c{c}\~ao.
\end{itemize}


Conclui-se que a plataforma vest\'ivel modular multissensor foi implementada e integrada com os mecanismos arquiteturais propostos. A continuidade do trabalho deve priorizar a caracteriza\c{c}\~ao quantitativa desses mecanismos e, em seguida, as expans\~oes de energia, armazenamento, conectividade, localiza\c{c}\~ao e algoritmos.

\end{document}
